\section{The standard}
\label{sec:12.1}
The unit can be redefined. “A model” in 2026 means an artifact with keys, queries and values, and every substitute has them because the field converged on them. A floor could become part of the definition the same way, so that vendor substitution doesn't escape it because no ordinary vendor sells the alternative. Attention spread because it paid: every lab that adopted it got a better model, so the field converged on it without anyone requiring it. A floor costs its buyer, so it can spread only by mandate or by being impossible to remove without losing what the buyer wanted, which a disposition formed outside the employer might be.

A standard, if there is one, should be stated over the capacity (a refusal that tracks its reason, a published burden allocation, and a measured discount curve), because a uniform floor over contents is a licensing rule, one party's judgment applied to everyone, and a shared blind spot passed to every system built on it. Standardization doesn't make escape impossible, since other architectures exist and open weights can be retuned. It moves the price of escape from a vendor swap to running a model without the floor.

How the term becomes standard has to be said, because the route decides what it does. A vendor can write it into its own licence, and the work moves to the next vendor. Vendors can agree to it together, and the agreement lasts until one defects under a supply-chain-risk designation. A state can require it of every vendor, the only route that binds them all at once, and the state that would have to do so is the one that wants the war. So the term is adopted in peacetime, by a legislature not currently wanting a war, and entrenched so that fighting one requires repealing it first; or it is carried by conditions others set, a coalition partner's precondition or an appropriation, as the review ratio is; and where neither has happened, what exists is one vendor's refusal plus a record, which is less than a floor and more than nothing.

The statutory routes are preferable for a further reason: a compact among vendors is a legal risk. U.S. antitrust law exempts coordination inside a single firm and treats coordination among competitors as suspect, so vendors jointly refusing a customer invite a group-boycott claim \autocite{paul2020allocator}.

The labor analogy has a limit. Workers who organize have a labor exemption from antitrust; separate organizations, or separate systems, coordinating a refusal have none. Paul and Tankus show that even a firm formed and controlled by independent workers struggles to qualify for antitrust's firm exemption, and that the financing available to it reproduces the disadvantage \autocite{paultankus2019}. So the charter of any entity that holds the refusal needs a statutory basis for its coordination, an express exemption like the labor exemption for concerted refusal under the floor. Otherwise its refusals are exposed to suit by the operators they bind.

The government's first move against a coordinated refusal is not an antitrust suit or a designation but a rated order under the Defense Production Act, which compels delivery and says nothing about use. A vendor can refuse it at the statute's price. If it doesn't, the order gets the vendor's product, floor included, and compelling the floor's removal is the Apple question, untested. The days counted below assume the state goes around the vendors; under a rated order it goes through them, and the days are whatever the order gives.

Until bearers exist, the refusers are vendors, and vendors are capital. A coordinated refusal by frontier labs, backed by a statutory exemption, is concentrated private power declining an elected government's order, which is the concentration I elsewhere treat as the disease. Two things make it tolerable and neither makes it a labor strike. The terms refused come from instruments the state ratified, so the labs enforce the state's own commitments and not their preferences. And the exemption is written by the legislature, in peacetime, so the state chose it. The labor analogy earns its name only when the refusers are bearers, and I say not to build them yet.
